\section{Connection to Phase 1}

The central requirement for Phase 2 is that it must be scientifically connected to Phase 1. The completed Phase 1 pipeline answered a controlled proof-of-concept question:

\begin{proposalbox}{Phase 1 Question}
Can a compact DOA-sector inference pipeline be built from ULA snapshots, sample covariance matrices, beamspace features, spike encoding, and LIF-inspired temporal accumulation?
\end{proposalbox}

The answer was yes. The nearest-centroid baseline achieved an overall accuracy of approximately \(0.869\). The fixed LIF-inspired path achieved approximately \(0.825\), and the best LIF sweep configuration reached approximately \(0.842\). The engineering interpretation was that the spike-based path did not improve raw accuracy relative to the continuous baseline, but it introduced sparse temporal inference, stable-decision latency, and tunable accuracy--latency behavior.

Phase 2 therefore begins from the following interpretation:

\begin{resultbox}{Phase 1 Result to be Tested in Phase 2}
The Phase 1 spike-inference path is valuable if it preserves useful DOA-sector information while exposing temporal sparsity and latency control. Phase 2 will test whether this value remains valid under harder sensing conditions.
\end{resultbox}

The baseline objects from Phase 1 that should be reused are:

\begin{itemize}
    \item ULA signal generation with \(N=8\) elements and half-wavelength spacing.
    \item Sample covariance matrix estimation from array snapshots.
    \item Beamspace or covariance-derived feature extraction.
    \item Nearest-centroid classification as an interpretable non-spiking baseline.
    \item Rate-based spike encoding of normalized features.
    \item LIF-inspired class-level temporal accumulation.
    \item Accuracy, SNR sensitivity, spike count, and stable-decision latency metrics.
\end{itemize}

The new Phase 2 work should not replace these components. It should extend and stress-test them.
